% titel: Schwerpunkt: Person geht auf ruhendem Boot
% skript: Kap. Schwerpunkt, Impulserhaltung

Eine Person bewegt sich innerhalb eines abgeschlossenen Systems (hier: auf
einem Boot, das ohne äussere Reibung auf dem Wasser liegt). Da im System
keine horizontalen äusseren Kräfte wirken, bleibt der gemeinsame Schwerpunkt
von Person und Boot am selben Ort. Aus dieser einzigen Tatsache folgt, wie
weit sich das Boot bewegt.

\begin{werkzeugkasten}[Werkzeuge]
\textbf{Schwerpunkt} eines Systems aus zwei Massen $m$ und $M$ an den Orten
$x_m$ und $x_M$:
\begin{equation}
  x_S = \frac{m\,x_m + M\,x_M}{m+M}\;.
\end{equation}

\textbf{Ohne äussere Kräfte} gilt $F_\mathrm{ext}=0$, der Gesamtimpuls $p$
ist konstant; der Schwerpunkt bleibt entweder in Ruhe oder in gleichförmiger
Bewegung. Verschiebt sich die eine Masse, so verschiebt sich die andere so,
dass $x_S$ unverändert bleibt.

\textbf{Relativbewegung:} Bewegt sich die Person um $\Delta x$ \emph{relativ
zum Boot}, so ist ihre Bewegung relativ zum Wasser die Summe aus
Bootgeschwindigkeit und relativer Geschwindigkeit. Das wird bei der
Momentum-Betrachtung gebraucht.
\end{werkzeugkasten}

\begin{werkzeugkasten}[Vorgehen in allgemeiner Form]
Schwerpunkt zu Beginn und am Ende ansetzen und gleichsetzen; daraus den
Versatz des Bootes bestimmen. Für Teil (b) den Impulssatz
$p = m(u+v)+Mv = 0$ mit der Relativgeschwindigkeit $u$ aufstellen.
\end{werkzeugkasten}

\begin{achtung}[Typische Fallen]
\begin{itemize}
  \item Nur die Verschiebung \emph{relativ zum Boot} ist die grosse Strecke
       $L$; die Person bewegt sich relativ zum Wasser anders.
  \item Vorzeichen: Person und Boot bewegen sich in entgegengesetzte
       Richtungen relativ zum Wasser.
\end{itemize}
\end{achtung}
